\chapter{Lewisstructuren, resonantie en VSEPR}\label{ch:b1:lewis-resonance-vsepr}

Het ozonmolecuul, \ce{O3}, bestaat uit drie zuurstofatomen in een
gebogen keten. Teken het met de regels van Boek 1 en de ene binding komt
dubbel uit, de andere enkel: de ene binding zou kort moeten zijn, de
andere lang. Toch vindt microgolfspectroscopie de twee \ce{O-O}-bindingen
van ozon precies even lang, \qty{127.8}{pm}, tussen de dubbele binding
van \ce{O2} (\qty{120.8}{pm}) en de enkele binding van waterstofperoxide
(\qty{147.5}{pm}) in. Geen enkele tekening klopt; twee tekeningen samen
wel. Dit hoofdstuk maakt het lewismodel nauwkeurig --- \omterm{def:b1:lewis-resonance-vsepr:formal-charge}{formele ladingen},
octetten die onvolledig of overschreden zijn, resonantie tussen meerdere
structuren --- en gebruikt het daarna om de vorm van moleculen te
voorspellen en de \omterm{def:b1:lewis-resonance-vsepr:dipole}{dipoolmomenten} die daaruit volgen.
% ledger: geo:O3.rOO, re:O2, geo:H2O2.rOO


\begin{recall}
Boek 1 (jaar 10) beschreef de \omterm{def:b1:lewis-resonance-vsepr:lewis}{covalente binding} als een gedeeld
elektronenpaar, tekende lewisstructuren met bindende en \omterm{def:b1:lewis-resonance-vsepr:lewis}{vrije
elektronenparen}, en voorspelde vier vormen (lineair, gebogen, trigonaal,
tetraëdrisch). De \omterm{def:b1:periodicity:electronegativity}{elektronegativiteit} en de verschillen ervan werden in
\cref{ch:b1:periodicity} kwantitatief gemaakt; de \omterm{def:b1:quantum-numbers:configuration}{valentie-elektronen} van
een atoom lees je af uit zijn configuratie
(\cref{def:b1:quantum-numbers:configuration}).
\end{recall}

\section{Lewisstructuren en formele ladingen}

\begin{definition}[Covalente binding, lewisstructuur]\label{def:b1:lewis-resonance-vsepr:lewis}
Een \emph{covalente binding}\index{covalente binding} is een
elektronenpaar dat door twee atomen gedeeld wordt, het \emph{bindend
elektronenpaar}\index{bindend elektronenpaar}; twee of drie paren vormen
een dubbele of een drievoudige binding. Een paar \omterm{def:b1:quantum-numbers:configuration}{valentie-elektronen} dat
bij één enkel atoom hoort, is een \emph{vrij
elektronenpaar}\index{vrij elektronenpaar}. Een
\emph{lewisstructuur}\index{lewisstructuur} toont elk
valentie-elektron van een molecuul of ion, als bindingen (streepjes
tussen atomen) en vrije elektronenparen (streepjes of paren punten op een
atoom). Volgens de \emph{octetregel}\index{octetregel} zijn de atomen van
\omterm{def:b1:periodicity:table}{periode} 2 (C, N, O, F) in een stabiele structuur omringd door vier paren,
acht elektronen; volgens de \emph{duetregel}\index{duetregel} is
waterstof omringd door één paar.
\end{definition}

\begin{definition}[Formele lading]\label{def:b1:lewis-resonance-vsepr:formal-charge}
De \emph{formele lading}\index{formele lading} van een atoom in een
\omterm{def:b1:lewis-resonance-vsepr:lewis}{lewisstructuur} is
\[
  q_F = N_v - N_{\text{vrij}} - \tfrac12 N_{\text{bind}},
\]
waarin $N_v$ het aantal \omterm{def:b1:quantum-numbers:configuration}{valentie-elektronen} van het vrije atoom is,
$N_{\text{vrij}}$ het aantal elektronen in zijn \omterm{def:b1:lewis-resonance-vsepr:lewis}{vrije elektronenparen} en
$N_{\text{bind}}$ het aantal elektronen in de bindingen die het vormt. Ze
wordt als $\oplus$ of $\ominus$ naast het atoom geschreven wanneer ze niet
nul is.
\end{definition}

\begin{proposition}[De formele ladingen tellen op tot de lading]\label{prop:b1:lewis-resonance-vsepr:formal-sum}
De som van de \omterm{def:b1:lewis-resonance-vsepr:formal-charge}{formele ladingen} van de atomen van een \omterm{def:b1:lewis-resonance-vsepr:lewis}{lewisstructuur} is
gelijk aan de totale lading van het molecuul of ion.
\end{proposition}

\begin{proof}
De som van $q_F$ over de atomen is $\sum N_v - (\text{elektronen in vrije
paren} + \text{bindingselektronen})$, omdat de elektronen van elke
binding voor de helft bij elk van haar twee atomen worden geteld. De
uitdrukking tussen haakjes is het totale aantal \omterm{def:b1:quantum-numbers:configuration}{valentie-elektronen} dat in
de structuur getekend is, en dat is gelijk aan $\sum N_v$ min de lading
$z$ van het deeltje. Dus is $\sum q_F = z$.
\end{proof}

\begin{method}[Een lewisstructuur tekenen]\label{met:b1:lewis-resonance-vsepr:lewis}
Zo teken je de \omterm{def:b1:lewis-resonance-vsepr:lewis}{lewisstructuur} van een molecuul of ion:
\begin{enumerate}
  \item tel de \omterm{def:b1:quantum-numbers:configuration}{valentie-elektronen} $N = \sum N_v - z$ en de paren $N/2$;
  \item teken het skelet: het minst elektronegatieve atoom (niet H) staat
        meestal centraal; waterstof en fluor zijn altijd eindstandig;
  \item vul de octetten van de eindstandige atomen aan met \omterm{def:b1:lewis-resonance-vsepr:lewis}{vrije paren} en
        plaats daarna de overige paren op het centrale atoom;
  \item ontbreekt het centrale atoom een octet, maak dan van \omterm{def:b1:lewis-resonance-vsepr:lewis}{vrije paren}
        van zijn buren meervoudige bindingen;
  \item bereken de \omterm{def:b1:lewis-resonance-vsepr:formal-charge}{formele ladingen}; kies uit de mogelijke structuren die
        met de minste \omterm{def:b1:lewis-resonance-vsepr:formal-charge}{formele ladingen}, en met de negatieve ladingen op
        de meest elektronegatieve atomen.
\end{enumerate}
\end{method}

\begin{example}[Salpeterzuur]\label{ex:b1:lewis-resonance-vsepr:hno3}
\ce{HNO3}: $N = 1 + 5 + 3 \times 6 = 24$ elektronen, 12 paren. Skelet:
stikstof centraal, gebonden aan drie zuurstofatomen, waarvan er één de
waterstof draagt. Worden de octetten met enkele bindingen aangevuld, dan
houdt stikstof zes elektronen over; een \omterm{def:b1:lewis-resonance-vsepr:lewis}{vrij paar} van een zuurstofatoom
wordt een \ce{N=O}-binding. \omterm{def:b1:lewis-resonance-vsepr:formal-charge}{Formele ladingen}: stikstof $5 - 0 - 4 = +1$;
het enkelvoudig gebonden zuurstofatoom zonder waterstof $6 - 6 - 1 = -1$;
de andere 0. Het totaal is 0, zoals
\cref{prop:b1:lewis-resonance-vsepr:formal-sum} vereist.
\end{example}

\begin{omfigure}
\setchemfig{atom sep=2.4em}
\begin{tikzpicture}[font=\small]
  \node at (0,0) {\large\chemfig{H-\charge{90=\|,270=\|}{O}-\charge{90:2pt=$\scriptstyle\oplus$}{N}(=[:45]\charge{0=\|,90=\|}{O})-[:-45]\charge{0=\|,240=\|,300=\|,45:3pt=$\scriptstyle\ominus$}{O}}};
  \node[below] at (0,-1.5) {salpeterzuur, \ce{HNO3}};
  \node at (5.2,0) {\large\chemfig{\charge{180=\|}{N}~\charge{0=\|}{N}}};
  \node[below] at (5.2,-1.5) {distikstof, \ce{N2}};
  \node at (9.6,0) {\large\chemfig{\charge{135=\|,225=\|,90:2pt=$\scriptstyle\ominus$}{N}=\charge{90:2pt=$\scriptstyle\oplus$}{N}=\charge{45=\|,315=\|}{O}}};
  \node[below] at (9.6,-1.5) {distikstofmonoxide, \ce{N2O}};
\end{tikzpicture}

{\small Lewisstructuren met \omterm{def:b1:lewis-resonance-vsepr:lewis}{vrije elektronenparen} en \omterm{def:b1:lewis-resonance-vsepr:formal-charge}{formele ladingen}.}
\end{omfigure}

\section{Voorbij het octet}

\begin{definition}[Elektronendeficiënte en hypervalente moleculen]\label{def:b1:lewis-resonance-vsepr:hypervalent}
Een molecuul waarin een atoom minder dan acht \omterm{def:b1:quantum-numbers:configuration}{valentie-elektronen} heeft,
is \emph{elektronendeficiënt}\index{elektronendeficient@elektronendeficiënt}:
in \ce{BF3} heeft boor er zes. Een molecuul waarin een atoom uit \omterm{def:b1:periodicity:table}{periode}
3 of verder omringd is door meer dan vier paren, is
\emph{hypervalent}\index{hypervalent}: in \ce{PCl5} vormt fosfor vijf
bindingen, in \ce{SF6} vormt zwavel er zes. Moleculen met een oneven
aantal elektronen, zoals \ce{NO} en \ce{NO2}, kunnen niet op elk atoom
aan de \omterm{def:b1:lewis-resonance-vsepr:lewis}{octetregel} voldoen: één elektron is ongepaard.
\end{definition}

\begin{remark}[Waarom periode 3 en niet periode 2]\label{rem:b1:lewis-resonance-vsepr:hypervalent}
Stikstof vormt \ce{NF3} maar nooit \ce{NF5}, terwijl fosfor zowel
\ce{PF3} als \ce{PF5} vormt. Een atoom uit \omterm{def:b1:periodicity:table}{periode} 2 is te klein om aan
meer dan vier buren te binden; een atoom uit \omterm{def:b1:periodicity:table}{periode} 3 is groot genoeg
voor vijf of zes. Tekeningen van \ce{SO4^{2-}} of \ce{H2SO4} met dubbele
bindingen aan zwavel (12 elektronen rond S) en met enkele bindingen en
\omterm{def:b1:lewis-resonance-vsepr:formal-charge}{formele ladingen} (een octet) komen allebei voor; de tweede geeft de
ladingsverdeling beter weer, en beide leveren de juiste geometrie.
\end{remark}

\begin{definition}[Lewiszuur, lewisbase, datieve binding]\label{def:b1:lewis-resonance-vsepr:lewis-acid}
Een \emph{lewiszuur}\index{lewiszuur} is een deeltje dat een
elektronenpaar kan opnemen op een vrije plaats in zijn valentieschil
(\ce{BF3}, \ce{AlCl3}, \ce{H+}, metaalkationen); een
\emph{lewisbase}\index{lewisbase} is een deeltje met een \omterm{def:b1:lewis-resonance-vsepr:lewis}{vrij
elektronenpaar} dat het kan delen (\ce{NH3}, \ce{H2O}, \ce{F-}). Wanneer de
base het paar aan het zuur geeft, is de gevormde binding een
\emph{datieve binding}\index{datieve binding}: een \omterm{def:b1:lewis-resonance-vsepr:lewis}{covalente binding}
waarvan beide elektronen van één partner afkomstig zijn.
\end{definition}

\begin{example}[Ammoniak en boortrifluoride]\label{ex:b1:lewis-resonance-vsepr:adduct}
\ce{NH3 + BF3 -> H3NBF3}. In het adduct voltooit boor zijn octet;
stikstof, dat nu zijn \omterm{def:b1:lewis-resonance-vsepr:lewis}{vrije paar} deelt, draagt de \omterm{def:b1:lewis-resonance-vsepr:formal-charge}{formele lading} $+1$ en
boor $-1$. Eenmaal gevormd is de \ce{N-B}-binding een gewone \omterm{def:b1:lewis-resonance-vsepr:lewis}{covalente
binding}.
\end{example}

\begin{omfigure}
\setchemfig{atom sep=2.1em}
\begin{tikzpicture}[font=\small]
  \node at (0,0) {\large\chemfig{H-\charge{90=\|}{N}(-[:240]H)(-[:300]H)}};
  \node at (1.6,0) {$+$};
  \node at (3.1,0) {\large\chemfig{B(-[:90]F)(-[:210]F)(-[:330]F)}};
  \draw[-{Stealth}, thick] (4.3,0) -- (5.3,0);
  \node at (7.4,0) {\large\chemfig{H-\charge{45:1pt=$\scriptstyle\oplus$}{N}(-[:90]H)(-[:270]H)-\charge{45:1pt=$\scriptstyle\ominus$}{B}(-[:90]F)(-[:270]F)-F}};
\end{tikzpicture}

{\small Een \omterm{def:b1:lewis-resonance-vsepr:lewis-acid}{datieve binding}: het vrije elektronenpaar van ammoniak (een
\omterm{def:b1:lewis-resonance-vsepr:lewis-acid}{lewisbase}) vult de vrije plaats van boortrifluoride (een \omterm{def:b1:lewis-resonance-vsepr:lewis-acid}{lewiszuur}). De
\omterm{def:b1:lewis-resonance-vsepr:lewis}{vrije paren} van fluor zijn niet getekend.}
\end{omfigure}

\section{Resonantie}

\begin{definition}[Resonantiestructuren, resonantiehybride]\label{def:b1:lewis-resonance-vsepr:resonance}
Wanneer de elektronen van een molecuul over meerdere lewisstructuren
verdeeld kunnen worden die alleen verschillen in de plaats van de
$\pi$-bindingen en de \omterm{def:b1:lewis-resonance-vsepr:lewis}{vrije paren} --- terwijl de kernen op hun plaats
blijven ---, is elk daarvan een
\emph{resonantiestructuur}\index{resonantiestructuur} (of mesomere
grensstructuur), met de andere verbonden door een dubbele pijl
$\leftrightarrow$. Het echte molecuul is geen van alle: het is de
\emph{resonantiehybride}\index{resonantiehybride}, één enkele structuur
waarvan de elektronenverdeling een gewogen gemiddelde is van die van de
grensstructuren.
\end{definition}

\begin{remark}[Een pijl, geen evenwicht]\label{rem:b1:lewis-resonance-vsepr:arrow}
De pijl $\leftrightarrow$ beschrijft geen reactie: ozon springt niet van
de ene structuur naar de andere. De hybride is één molecuul, beschreven
met meerdere tekeningen omdat één tekening niet volstaat. Gebruik nooit
$\rightleftharpoons$ tussen \omterm{def:b1:lewis-resonance-vsepr:resonance}{resonantiestructuren}.
\end{remark}

\begin{method}[Resonantiestructuren opschrijven en wegen]\label{met:b1:lewis-resonance-vsepr:resonance}
Zo vind je de \omterm{def:b1:lewis-resonance-vsepr:resonance}{resonantiestructuren} van een deeltje:
\begin{enumerate}
  \item vertrek van één \omterm{def:b1:lewis-resonance-vsepr:lewis}{lewisstructuur}; zoek een \omterm{def:b1:lewis-resonance-vsepr:lewis}{vrij paar} of een
        $\pi$-binding naast een $\pi$-binding, een vrije plaats of een
        positieve lading;
  \item verplaats paren met gebogen pijlen (een \omterm{def:b1:lewis-resonance-vsepr:lewis}{vrij paar} wordt een
        $\pi$-binding, een $\pi$-binding wordt een \omterm{def:b1:lewis-resonance-vsepr:lewis}{vrij paar}), zonder ooit
        een kern te verplaatsen en zonder ooit het octet van een atoom
        uit \omterm{def:b1:periodicity:table}{periode} 2 te overschrijden;
  \item bereken de \omterm{def:b1:lewis-resonance-vsepr:formal-charge}{formele ladingen} opnieuw;
  \item weeg de structuren: de belangrijkste voldoen aan de \omterm{def:b1:lewis-resonance-vsepr:lewis}{octetregel},
        hebben de minste \omterm{def:b1:lewis-resonance-vsepr:formal-charge}{formele ladingen} en leggen de negatieve ladingen
        op de meest elektronegatieve atomen; gelijkwaardige structuren
        wegen even zwaar.
\end{enumerate}
\end{method}

\begin{example}[Ozon, nitraat, benzeen]\label{ex:b1:lewis-resonance-vsepr:examples}
Ozon heeft twee gelijkwaardige structuren: elke \ce{O-O}-binding is in de
ene dubbel en in de andere enkel, zodat beide een bindingsorde $1.5$ en
dezelfde lengte hebben, \qty{127.8}{pm}, tussen \ce{O=O} en \ce{O-O} in.
Het nitraation \ce{NO3-} heeft drie gelijkwaardige structuren; elke
\ce{N-O}-binding heeft orde $4/3$. Benzeen, \ce{C6H6}, heeft twee
kekulestructuren; zijn zes \ce{C-C}-bindingen zijn even lang,
\qty{139.7}{pm}, tussen de enkele binding van ethaan (\qty{153.6}{pm})
en de dubbele binding van etheen (\qty{133.9}{pm}) in.
% ledger: geo:O3.rOO, geo:C6H6.rCC, geo:C2H6.rCC, geo:C2H4.rCC
\end{example}

\begin{omfigure}
\vspace*{10pt}
\large
\setchemfig{atom sep=2.4em}
\schemestart
\chemfig{\charge{135=\|,225=\|}{O}=[:-30]\charge{270=\|,90:2pt=$\scriptstyle\oplus$}{O}-[:30]\charge{90=\|,0=\|,270=\|,45:5pt=$\scriptstyle\ominus$}{O}}
\arrow{<->}
\chemfig{\charge{90=\|,180=\|,270=\|,135:5pt=$\scriptstyle\ominus$}{O}-[:-30]\charge{270=\|,90:2pt=$\scriptstyle\oplus$}{O}=[:30]\charge{45=\|,315=\|}{O}}
\schemestop

\medskip
\schemestart
\chemfig{*6(=-=-=-)}
\arrow{<->}
\chemfig{*6(-=-=-=)}
\schemestop
\hspace{3em}
\chemfig{\charge{270:2pt=$\scriptstyle\oplus$}{N}(=[:90]\charge{45=\|,135=\|}{O})(-[:210]\charge{150=\|,240=\|,300=\|,90:2pt=$\scriptstyle\ominus$}{O})-[:330]\charge{30=\|,300=\|,240=\|,90:2pt=$\scriptstyle\ominus$}{O}}
\normalsize
\par\vspace{14pt}
{\small Resonantie in ozon (een \omterm{def:b1:lewis-resonance-vsepr:lewis}{vrij paar} van het rechtse zuurstofatoom
wordt een $\pi$-binding, terwijl de $\pi$-binding links een \omterm{def:b1:lewis-resonance-vsepr:lewis}{vrij paar}
wordt) en in benzeen, en een van de drie gelijkwaardige structuren van
het nitraation, waarvan de lading over de drie zuurstofatomen verdeeld
is.}
\end{omfigure}

\begin{definition}[$\sigma$-binding, $\pi$-binding]\label{def:b1:lewis-resonance-vsepr:sigma-pi}
Een enkele binding is een \emph{$\sigma$-binding}\index{sigma-binding@$\sigma$-binding}:
haar elektronenpaar ligt langs de as die de twee kernen verbindt, door de
frontale overlapping van twee orbitalen. De tweede en de derde binding
van een meervoudige binding zijn \emph{$\pi$-bindingen}\index{pi-binding@$\pi$-binding}:
hun paren liggen aan weerszijden van de as, door de zijdelingse
overlapping van twee p-orbitalen die er loodrecht op staan. Een dubbele
binding bestaat uit één $\sigma$- en één $\pi$-binding; een drievoudige
binding uit één $\sigma$- en twee $\pi$-bindingen.
\end{definition}
\begin{omfigure}
\begin{tikzpicture}[font=\small]
  % sigma: two lobes overlapping head-on
  \fill[omDef!35, opacity=0.8] (0,0) ellipse (0.9 and 0.42);
  \fill[omProp!35, opacity=0.8] (1.1,0) ellipse (0.9 and 0.42);
  \fill (0.1,0) circle (1.6pt); \fill (1.0,0) circle (1.6pt);
  \draw[dashed, black!50] (-1.2,0) -- (2.3,0);
  \node[below] at (0.55,-1.05) {$\sigma$: frontaal, op de as};
  % pi: two p orbitals side by side
  \begin{scope}[xshift=5.2cm]
    \foreach \x/\c in {0/omDef, 1.2/omProp} {
      \fill[\c!35, opacity=0.85] (\x,0.45) ellipse (0.32 and 0.45);
      \fill[\c!15, opacity=0.85] (\x,-0.45) ellipse (0.32 and 0.45);
      \draw[\c!70!black] (\x,0.45) ellipse (0.32 and 0.45) (\x,-0.45) ellipse (0.32 and 0.45);
      \fill (\x,0) circle (1.6pt);
    }
    \draw[dashed, black!50] (-0.7,0) -- (1.9,0);
    \draw[omThm, thick, densely dotted] (0.25,0.75) -- (0.95,0.75) (0.25,-0.75) -- (0.95,-0.75);
    \node[below] at (0.6,-1.05) {$\pi$: zijdelings, boven en onder};
  \end{scope}
\end{tikzpicture}

{\small De twee soorten overlapping. Een $\sigma$-binding laat rotatie om
haar as toe; een $\pi$-binding niet, en daarom is een dubbele binding
star (\cref{ch:b1:stereochemistry-in-depth}).}
\end{omfigure}

\section{Het VSEPR-model}

\begin{definition}[VSEPR-model, sterisch getal]\label{def:b1:lewis-resonance-vsepr:vsepr}
In het \emph{VSEPR-model}\index{VSEPR-model} (\textit{valence-shell
electron-pair repulsion}, afstoting van de elektronenparen in de
valentieschil) stoten de paren rond een centraal atoom A elkaar af en
nemen ze de schikking aan die hen zo ver mogelijk van elkaar houdt. Een
molecuul wordt geschreven als \ce{AX_nE_m}, waarin $n$ het aantal atomen
is dat aan A gebonden is (een meervoudige binding telt één keer) en $m$
het aantal \omterm{def:b1:lewis-resonance-vsepr:lewis}{vrije paren} van A. Het \emph{sterisch getal}\index{sterisch getal}
$n + m$ legt de schikking van de paren vast: 2 lineair, 3 trigonaal vlak,
4 tetraëdrisch, 5 trigonaal bipiramidaal, 6 octaëdrisch. De vorm van het
molecuul is alleen de schikking van zijn atomen.
\end{definition}

\begin{proposition}[VSEPR-geometrieën]\label{prop:b1:lewis-resonance-vsepr:geometries}
Het \omterm{def:b1:lewis-resonance-vsepr:vsepr}{VSEPR-model} voorspelt de vormen en ideale hoeken van de tabel
hieronder; de gemeten hoeken liggen er dicht bij.
\begin{center}
\small
\begin{tabular}{llllc}
\hline
type & schikking van de paren & vorm & voorbeeld & gemeten hoek \\
\hline
\ce{AX2} & lineair & lineair & \ce{CO2} & $180^\circ$ \\
\ce{AX3} & trigonaal vlak & trigonaal vlak & \ce{BF3} & $120^\circ$ \\
\ce{AX2E} & trigonaal vlak & gebogen & \ce{SO2} & $119.5^\circ$ \\
\ce{AX4} & tetraëdrisch & tetraëdrisch & \ce{CH4} & $109.5^\circ$ \\
\ce{AX3E} & tetraëdrisch & trigonaal piramidaal & \ce{NH3} & $106.7^\circ$ \\
\ce{AX2E2} & tetraëdrisch & gebogen & \ce{H2O} & $104.5^\circ$ \\
\ce{AX5} & trigonaal bipiramidaal & trigonaal bipiramidaal & \ce{PCl5} & $90^\circ$, $120^\circ$ \\
\ce{AX4E} & trigonaal bipiramidaal & wipplank & \ce{SF4} & $101.6^\circ$, $173.1^\circ$ \\
\ce{AX3E2} & trigonaal bipiramidaal & T-vormig & \ce{ClF3} & $87.5^\circ$ \\
\ce{AX2E3} & trigonaal bipiramidaal & lineair & \ce{XeF2} & $180^\circ$ \\
\ce{AX6} & octaëdrisch & octaëdrisch & \ce{SF6} & $90^\circ$ \\
\ce{AX5E} & octaëdrisch & vierkant piramidaal & \ce{BrF5} & --- \\
\ce{AX4E2} & octaëdrisch & vierkant vlak & \ce{XeF4} & --- \\
\hline
\end{tabular}
\end{center}
% ledger: geo:CO2.aOCO, geo:BF3.aFBF, geo:SO2.aOSO, geo:CH4.aHCH,
% geo:NH3.aHNH, geo:H2O.aHOH, geo:SF6.aFSF, geo:SF4.aFSF2, geo:SF4.aFSF,
% geo:ClF3.aFClF, geo:XeF2.aFXeF
\end{proposition}
\admitted


\begin{proposition}[Vrije paren verkleinen de hoeken]\label{prop:b1:lewis-resonance-vsepr:lone-pair-angles}
Een \omterm{def:b1:lewis-resonance-vsepr:lewis}{vrij paar}, dat door één kern wordt vastgehouden, neemt meer ruimte in
dan een \omterm{def:b1:lewis-resonance-vsepr:lewis}{bindend paar} en stoot sterker af: de afstotingen nemen af in de
volgorde \omterm{def:b1:lewis-resonance-vsepr:lewis}{vrij paar}/\omterm{def:b1:lewis-resonance-vsepr:lewis}{vrij paar} $>$ \omterm{def:b1:lewis-resonance-vsepr:lewis}{vrij paar}/\omterm{def:b1:lewis-resonance-vsepr:lewis}{bindend paar} $>$ \omterm{def:b1:lewis-resonance-vsepr:lewis}{bindend
paar}/\omterm{def:b1:lewis-resonance-vsepr:lewis}{bindend paar}. Daardoor worden de hoeken tussen de bindingen kleiner
naarmate er \omterm{def:b1:lewis-resonance-vsepr:lewis}{vrije paren} bijkomen ($109.5^\circ$ in \ce{CH4},
$106.7^\circ$ in \ce{NH3}, $104.5^\circ$ in \ce{H2O}), en in een
trigonale bipiramide nemen de \omterm{def:b1:lewis-resonance-vsepr:lewis}{vrije paren} de equatoriale posities in,
waar ze maar twee buren op $90^\circ$ hebben.
% ledger: geo:CH4.aHCH, geo:NH3.aHNH, geo:H2O.aHOH
\end{proposition}
\admitted

\begin{omfigure}
\setchemfig{atom sep=2.6em}
\renewcommand{\arraystretch}{1.2}
\begin{tabular}{ccccc}
\chemfig{O=C=O} &
\chemfig{B(-[:90]F)(-[:210]F)(-[:330]F)} &
\chemfig{C(-[:90]H)(-[:200]H)(<[:280]H)(<:[:335]H)} &
\chemfig{P(-[:90]Cl)(-[:270]Cl)(-[:180]Cl)(<[:320]Cl)(<:[:30]Cl)} &
\chemfig{S(-[:90]F)(-[:270]F)(-[:180]F)(-[:0]F)(<[:225]F)(<:[:45]F)} \\
\small \ce{AX2}, \ce{CO2} & \small \ce{AX3}, \ce{BF3} & \small \ce{AX4}, \ce{CH4} &
\small \ce{AX5}, \ce{PCl5} & \small \ce{AX6}, \ce{SF6} \\[16pt]
\chemfig{\charge{90=\:}{S}(=[:210]O)(=[:330]O)} &
\chemfig{\charge{90=\:}{N}(-[:200]H)(<[:280]H)(<:[:335]H)} &
\chemfig{\charge{60=\:,120=\:}{O}(-[:215]H)(-[:325]H)} &
\chemfig{\charge{0=\:}{S}(-[:90]F)(-[:270]F)(<[:210]F)(<:[:150]F)} &
\chemfig{Xe(-[:0]F)(-[:90]F)(-[:180]F)(-[:270]F)} \\
\small \ce{AX2E}, \ce{SO2} & \small \ce{AX3E}, \ce{NH3} & \small \ce{AX2E2}, \ce{H2O} &
\small \ce{AX4E}, \ce{SF4} & \small \ce{AX4E2}, \ce{XeF4} \\
\end{tabular}

\medskip
{\small Vormen voorspeld door VSEPR. Een wig wijst naar de lezer toe,
een gearceerde wig van de lezer af; gewone bindingen liggen in het vlak
van het papier. \omterm{def:b1:lewis-resonance-vsepr:lewis}{Vrije paren} van het centrale atoom zijn als paren punten
weergegeven; \ce{XeF4} is van voren getekend, met zijn twee \omterm{def:b1:lewis-resonance-vsepr:lewis}{vrije paren}
naar de lezer toe en van de lezer af gericht.}
\end{omfigure}

\begin{method}[Een vorm voorspellen]\label{met:b1:lewis-resonance-vsepr:vsepr}
Zo voorspel je de vorm rond een centraal atoom:
\begin{enumerate}
  \item teken de \omterm{def:b1:lewis-resonance-vsepr:lewis}{lewisstructuur} (\cref{met:b1:lewis-resonance-vsepr:lewis});
  \item tel de atomen die aan het centrale atoom gebonden zijn ($n$) en
        zijn \omterm{def:b1:lewis-resonance-vsepr:lewis}{vrije paren} ($m$); schrijf \ce{AX_nE_m};
  \item lees de schikking af uit $n + m$ en de vorm uit de posities van
        de $n$ atomen;
  \item corrigeer de ideale hoeken: \omterm{def:b1:lewis-resonance-vsepr:lewis}{vrije paren} en meervoudige bindingen
        nemen meer ruimte in dan enkele bindingen.
\end{enumerate}
\end{method}

\begin{definition}[Hybridisatie]\label{def:b1:lewis-resonance-vsepr:hybridisation}
De \emph{hybridisatie}\index{hybridisatie} van een atoom is een
beschrijving van zijn bindingen met orbitalen die in de door VSEPR
voorspelde richtingen wijzen: een atoom met \omterm{def:b1:lewis-resonance-vsepr:vsepr}{sterisch getal} 4 heet
sp$^3$ (vier gelijkwaardige orbitalen op $109.5^\circ$), met \omterm{def:b1:lewis-resonance-vsepr:vsepr}{sterisch
getal} 3 sp$^2$ (drie orbitalen op $120^\circ$ in een vlak, waarbij één
p-orbitaal loodrecht daarop overblijft), met \omterm{def:b1:lewis-resonance-vsepr:vsepr}{sterisch getal} 2 sp (twee
orbitalen op $180^\circ$, waarbij twee p-orbitalen overblijven). De
overgebleven p-orbitalen vormen de $\pi$-bindingen.
\end{definition}

\begin{example}[De koolstofatomen van de organische chemie]\label{ex:b1:lewis-resonance-vsepr:carbon}
In ethaan is elk koolstofatoom sp$^3$, tetraëdrisch; in etheen sp$^2$,
met de zes atomen in één vlak en hoeken van ongeveer $120^\circ$ (gemeten
\ce{H-C-H} $117.6^\circ$); in ethyn sp, met de vier atomen op één lijn.
De bindingslengten krimpen van 153.6 via 133.9 tot \qty{120.3}{pm}
naarmate er $\pi$-bindingen bijkomen.
% ledger: geo:C2H4.aHCH, geo:C2H6.rCC, geo:C2H4.rCC, geo:C2H2.rCC
\end{example}


\begin{remark}[Een beschrijving, geen oorzaak]\label{rem:b1:lewis-resonance-vsepr:hybrid-limit}
\omterm{def:b1:lewis-resonance-vsepr:hybridisation}{Hybridisatie} beschrijft een geometrie die al bekend is; ze verklaart
haar niet, en ze schiet tekort voor de energieën van de elektronen, die
de moleculaireorbitaaltheorie, in het deel van jaar 2, wel verklaart. Ze
blijft de dagelijkse taal van de organische chemie: een
``sp$^2$-koolstofatoom'' betekent een vlak koolstofatoom met een
p-orbitaal die vrij is voor een $\pi$-binding.
\end{remark}

\section{Polaire moleculen: dipoolmomenten}

\begin{definition}[Bindingsdipool, dipoolmoment]\label{def:b1:lewis-resonance-vsepr:dipole}
In een binding tussen atomen met een verschillende \omterm{def:b1:periodicity:electronegativity}{elektronegativiteit}
hellen de elektronen over naar het elektronegatievere atoom, dat een
partiële lading $-\delta e$ draagt, terwijl zijn partner $+\delta e$
draagt ($0 < \delta < 1$). De \emph{bindingsdipool}\index{bindingsdipool}
is de vector $\vec\mu = \delta e\,
\vec d$, met norm $\delta e d$, gericht van de negatieve naar de positieve
partiële lading, waarbij $\vec d$ de vector tussen beide is. Het
\emph{dipoolmoment}\index{dipoolmoment} van een molecuul is de
vectorsom van zijn bindingsdipolen (en van de bijdragen van zijn \omterm{def:b1:lewis-resonance-vsepr:lewis}{vrije
paren}); het wordt uitgedrukt in coulombmeter of in debye, $\qty{1}{\debye} =
\qty{3.336e-30}{C.m}$. Een molecuul met een dipoolmoment verschillend van
nul is een \emph{polair molecuul}\index{polair molecuul}.
% ledger: const:c (1 D = 1e-21/c C.m)
\end{definition}


\begin{remark}[Richtingsafspraak]\label{rem:b1:lewis-resonance-vsepr:convention}
De natuurkunde tekent de dipoolvector van de negatieve naar de positieve
lading; dat is de afspraak die hier gebruikt wordt. Oudere
scheikundeboeken tekenen een pijl met een doorgestreepte staart die naar
het negatieve uiteinde wijst; alleen de richting verschilt.
\end{remark}

\begin{proposition}[Dipool van een symmetrisch molecuul]\label{prop:b1:lewis-resonance-vsepr:dipole-sum}
Een molecuul \ce{AX_n} zonder \omterm{def:b1:lewis-resonance-vsepr:lewis}{vrije paren} op A (lineair \ce{AX2},
trigonaal \ce{AX3}, tetraëdrisch \ce{AX4}, \dots) heeft een \omterm{def:b1:lewis-resonance-vsepr:dipole}{dipoolmoment}
nul, hoe polair zijn bindingen ook zijn. Een gebogen molecuul \ce{AX2}
met hoek $\theta$, waarvan de bindingen de dipool $\mu_b$ hebben, heeft
het \omterm{def:b1:lewis-resonance-vsepr:dipole}{dipoolmoment}
\[
  \mu = 2\mu_b \cos\frac{\theta}{2} .
\]
\end{proposition}


\begin{proof}
In de symmetrische vormen zijn de bindingsvectoren even groot en zijn
hun richtingen symmetrisch rond A geschikt, zodat ze opgeteld nul geven
(voor lineair \ce{AX2} zijn ze tegengesteld; voor \ce{AX3} op
$120^\circ$ is de som van drie vectoren met dezelfde norm nul; voor
\ce{AX4} is de som invariant onder de rotaties van de tetraëder, en dus
nul). Voor het gebogen molecuul maakt elke bindingsvector de hoek
$\theta/2$ met de bissectrice; de componenten loodrecht op de
bissectrice heffen elkaar op en die erlangs tellen op:
$2\mu_b\cos(\theta/2)$.
\end{proof}
\begin{omfigure}
\begin{tikzpicture}[font=\small, >={Stealth}]
  % water
  \node[aO] (o) at (0,0) {};
  \node[aH] (h1) at (-0.95,-0.73) {};
  \node[aH] (h2) at (0.95,-0.73) {};
  \begin{scope}[on background layer]
    \draw[bond] (o.center) -- (h1.center); \draw[bond] (o.center) -- (h2.center);
  \end{scope}
  \draw[->, omProp, thick] (-0.25,0.15) -- (-0.95,-0.38);
  \draw[->, omProp, thick] (0.25,0.15) -- (0.95,-0.38);
  \draw[->, omThm, very thick] (0,-0.25) -- (0,-1.45);
  \node[omThm, right] at (0.05,-1.25) {$\vec\mu$};
  \node at (0,0.55) {$\delta^-$};
  \node at (-1.25,-1.05) {$\delta^+$}; \node at (1.25,-1.05) {$\delta^+$};
  \node[below] at (0,-1.7) {\ce{H2O}: $\mu = \qty{1.86}{\debye}$};
  % carbon dioxide
  \begin{scope}[xshift=5cm]
    \node[aO] (a) at (-1.1,0) {}; \node[aC] (c) at (0,0) {}; \node[aO] (b) at (1.1,0) {};
    \begin{scope}[on background layer]
      \foreach \d in {-1.6,1.6} {
        \draw[bond, line width=1.1pt] ([yshift=\d pt]a.center) -- ([yshift=\d pt]c.center);
        \draw[bond, line width=1.1pt] ([yshift=\d pt]c.center) -- ([yshift=\d pt]b.center);}
    \end{scope}
    \draw[->, omProp, thick] (-0.75,0.45) -- (-0.2,0.45);
    \draw[->, omProp, thick] (0.75,0.45) -- (0.2,0.45);
    \node[below] at (0,-1.7) {\ce{CO2}: bindingsdipolen heffen elkaar op, $\mu = 0$};
  \end{scope}
\end{tikzpicture}

{\small \omterm{def:b1:lewis-resonance-vsepr:dipole}{Bindingsdipolen} (oranje, van de negatieve naar de positieve
partiële lading) en hun som (rood). In het gebogen watermolecuul tellen
ze op; in het lineaire koolstofdioxide heffen ze elkaar op.}
% ledger: dip:H2O
\end{omfigure}


\begin{example}[Polaire en apolaire moleculen]\label{ex:b1:lewis-resonance-vsepr:dipoles}
\ce{CO2}, \ce{BF3}, \ce{CH4} en \ce{CCl4} hebben een \omterm{def:b1:lewis-resonance-vsepr:dipole}{dipoolmoment} nul.
\ce{H2O} (\qty{1.86}{\debye}), \ce{NH3} (\qty{1.48}{\debye}) en
\ce{SO2} (\qty{1.63}{\debye}) zijn polair. Van \ce{CH3Cl} over
\ce{CH2Cl2} naar \ce{CHCl3} daalt het moment, 1.87, 1.62,
\qty{1.04}{\debye}, om in \ce{CCl4} te verdwijnen: de
\ce{C-Cl}-bindingsdipolen heffen elkaar steeds meer op. Bij de
waterstofhalogeniden volgt het het verschil in \omterm{def:b1:periodicity:electronegativity}{elektronegativiteit}:
HF 1.83, HCl 1.09, HBr 0.83, HI \qty{0.45}{\debye}.
% ledger: dip:H2O, dip:NH3, dip:SO2, dip:CH3Cl, dip:CH2Cl2, dip:CHCl3,
% dip:CCl4, dip:HF, dip:HCl, dip:HBr, dip:HI
\end{example}


\section{Oefeningen}

\begin{exercise}[$\star$]\label{exo:b1:lewis-resonance-vsepr:1}
Teken de lewisstructuren van \ce{H2O}, \ce{NH3}, \ce{CO2}, \ce{HCN},
\ce{CH2O} (methanal) en \ce{NH4+}, en geef de \omterm{def:b1:lewis-resonance-vsepr:formal-charge}{formele ladingen}.
\end{exercise}

\begin{exercise}[$\star$]\label{exo:b1:lewis-resonance-vsepr:2}
Geef het type \ce{AX_nE_m}, de vorm en de benaderde hoek van \ce{H2S},
\ce{PCl3}, \ce{BF3}, \ce{SiH4} en \ce{CO3^{2-}} (in elke formule staat
het centrale atoom voorop).
\end{exercise}

\begin{exercise}[$\star$]\label{exo:b1:lewis-resonance-vsepr:3}
Welke van deze moleculen zijn polair: \ce{CCl4}, \ce{CH2Cl2}, \ce{BF3},
\ce{NH3}, \ce{CO2}, \ce{SO2}? Verantwoord vanuit de vormen.
\end{exercise}

\begin{exercise}[$\star$]\label{exo:b1:lewis-resonance-vsepr:4}
Wijs het \omterm{def:b1:lewis-resonance-vsepr:lewis-acid}{lewiszuur} en de \omterm{def:b1:lewis-resonance-vsepr:lewis-acid}{lewisbase} aan in \ce{H+ + H2O -> H3O+} en in
\ce{Cu^{2+} + 4NH3 -> [Cu(NH3)4]^{2+}}. Welke bindingen zijn datief?
\end{exercise}

\begin{exercise}[$\star\star$]\label{exo:b1:lewis-resonance-vsepr:5}
Schrijf twee \omterm{def:b1:lewis-resonance-vsepr:resonance}{resonantiestructuren} van het ethanoaation \ce{CH3COO-}. Wat
voorspellen ze voor de lengten van de twee \ce{C-O}-bindingen? Vergelijk
met methaanzuur \ce{HCOOH}, waarvan de twee \ce{C-O}-bindingen 120.2 en
\qty{134.3}{pm} meten.
% ledger: geo:H2CO2.rCO, geo:H2CO2.rCO2
\end{exercise}


\begin{exercise}[$\star\star$]\label{exo:b1:lewis-resonance-vsepr:6}
Teken de \omterm{def:b1:lewis-resonance-vsepr:lewis}{lewisstructuur} van het sulfaation \ce{SO4^{2-}} met vier
enkele \ce{S-O}-bindingen, bereken de \omterm{def:b1:lewis-resonance-vsepr:formal-charge}{formele ladingen} en voorspel de
vorm. Hoeveel \omterm{def:b1:lewis-resonance-vsepr:resonance}{resonantiestructuren} met twee \ce{S=O}-bindingen en zonder
lading op zwavel kun je opschrijven?
\end{exercise}

\begin{exercise}[$\star\star$]\label{exo:b1:lewis-resonance-vsepr:7}
Voorspel met VSEPR de vormen van \ce{SF4}, \ce{ClF3}, \ce{XeF2} en
\ce{XeF4}. Leg uit waarom de \omterm{def:b1:lewis-resonance-vsepr:lewis}{vrije paren} van een trigonale bipiramide in
het equatoriale vlak liggen.
\end{exercise}

\begin{exercise}[$\star\star$]\label{exo:b1:lewis-resonance-vsepr:8}
De hoeken \ce{H-X-H} bedragen $104.5^\circ$ in \ce{H2O} en $92.1^\circ$
in \ce{H2S}; $106.7^\circ$ in \ce{NH3} en $93.3^\circ$ in \ce{PH3}. Stel
een verklaring voor op basis van de grootte en de \omterm{def:b1:periodicity:electronegativity}{elektronegativiteit}
van het centrale atoom.
% ledger: geo:H2O.aHOH, geo:H2S.aHSH, geo:NH3.aHNH, geo:PH3.aHPH
\end{exercise}


\begin{exercise}[$\star\star$]\label{exo:b1:lewis-resonance-vsepr:9}
Geef de \omterm{def:b1:lewis-resonance-vsepr:hybridisation}{hybridisatie} van elk koolstofatoom en van het stikstofatoom in
\ce{CH3-C#N} (ethaannitril) en in \ce{CH2=CH-CH3}, en het aantal
$\sigma$- en $\pi$-bindingen in elk molecuul.
\end{exercise}

\begin{exercise}[$\star\star\star$]\label{exo:b1:lewis-resonance-vsepr:10}
Het \omterm{def:b1:lewis-resonance-vsepr:dipole}{dipoolmoment} van \ce{NF3} is slechts \qty{0.24}{\debye}, tegen
\qty{1.48}{\debye} voor \ce{NH3}, hoewel de \ce{N-F}-binding polairder
is dan de \ce{N-H}-binding. Verklaar het verschil met behulp van de
\omterm{def:b1:lewis-resonance-vsepr:dipole}{bindingsdipolen} en het \omterm{def:b1:lewis-resonance-vsepr:lewis}{vrije paar} van stikstof.
% ledger: dip:NF3, dip:NH3
\end{exercise}


\begin{exercise}[$\star\star\star$]\label{exo:b1:lewis-resonance-vsepr:11}
Het \ce{SO2}-molecuul heeft een hoek van $119.5^\circ$ en een
\omterm{def:b1:lewis-resonance-vsepr:dipole}{dipoolmoment} van \qty{1.63}{\debye}. Bereken de dipool van één
\ce{S-O}-binding; leid daaruit, met de bindingslengte \qty{143.2}{pm},
de partiële lading $\delta$ op elk zuurstofatoom af.
% ledger: geo:SO2.aOSO, geo:SO2.rSO, dip:SO2
\end{exercise}


\begin{exercise}[$\star\star\star$]\label{exo:b1:lewis-resonance-vsepr:12}
Teken de belangrijkste lewisstructuren van distikstofmonoxide \ce{N2O}
(lineair, \ce{N-N-O}) en bereken de \omterm{def:b1:lewis-resonance-vsepr:formal-charge}{formele ladingen}. De gemeten
bindingslengten zijn \ce{N-N} \qty{112.8}{pm} en \ce{N-O}
\qty{118.4}{pm}; in \ce{N2} meet de drievoudige binding \qty{109.8}{pm}.
Welke structuur weegt het zwaarst?
% ledger: geo:N2O.rNN, geo:N2O.rNO, re:N2
\end{exercise}


\section{Probleem: de moleculen van een blikseminslag}

\begin{problem}[{Weekendprobleem --- de stikstof- en zuurstofdeeltjes die
ontstaan wanneer de bliksem lucht verhit: lewisstructuren, resonantie,
vormen, en de partiële ladingen van water}]
\label{pb:b1:lewis-resonance-vsepr:1}
Een bliksemschicht verhit de lucht tot duizenden graden: \ce{N2} en
\ce{O2} geven \ce{NO}, daarna \ce{NO2}, \ce{N2O}, ozon en ten slotte
salpeterzuur in de regen. Gegeven: $\qty{1}{\debye} = \qty{3.336e-30}{C.m}$,
$e =
\qty{1.602e-19}{C}$. Gemeten waarden worden gegeven waar nodig.

\medskip
\noindent\textbf{Deel I --- Lewisstructuren.}
\begin{enumerate}
  \item Tel de \omterm{def:b1:quantum-numbers:configuration}{valentie-elektronen} van \ce{N2}, \ce{NO}, \ce{NO2},
        \ce{N2O}, \ce{O3} en \ce{HNO3}.
  \item Teken de \omterm{def:b1:lewis-resonance-vsepr:lewis}{lewisstructuur} van \ce{N2}.
  \item Teken een \omterm{def:b1:lewis-resonance-vsepr:lewis}{lewisstructuur} van \ce{NO}. Waarom kan niet op beide
        atomen aan de \omterm{def:b1:lewis-resonance-vsepr:lewis}{octetregel} worden voldaan?
  \item Teken twee lewisstructuren van \ce{N2O}, met \ce{N=N=O} en met
        \ce{N#N-O}, en geef de \omterm{def:b1:lewis-resonance-vsepr:formal-charge}{formele ladingen}.
  \item Teken een \omterm{def:b1:lewis-resonance-vsepr:lewis}{lewisstructuur} van ozon en geef de \omterm{def:b1:lewis-resonance-vsepr:formal-charge}{formele ladingen}.
  \item Teken een \omterm{def:b1:lewis-resonance-vsepr:lewis}{lewisstructuur} van \ce{HNO3} en geef de \omterm{def:b1:lewis-resonance-vsepr:formal-charge}{formele ladingen}.
  \item Welke van de zes deeltjes hebben een \omterm{def:b1:quantum-numbers:unpaired}{ongepaard elektron}?
\end{enumerate}

\medskip
\noindent\textbf{Deel II --- Resonantie en bindingslengten.}
\begin{enumerate}[resume]
  \item Schrijf de twee \omterm{def:b1:lewis-resonance-vsepr:resonance}{resonantiestructuren} van ozon. Welke bindingsorde
        voorspellen ze?
  \item De \ce{O-O}-binding meet \qty{127.8}{pm} in ozon, \qty{120.8}{pm}
        in \ce{O2} en \qty{147.5}{pm} in \ce{H2O2}. Wordt de voorspelling
        bevestigd?
  \item Schrijf de drie \omterm{def:b1:lewis-resonance-vsepr:resonance}{resonantiestructuren} van \ce{NO3-} en geef de orde
        van elke \ce{N-O}-binding.
  \item In \ce{HNO3} meten de drie \ce{N-O}-bindingen 140.6, 121.1 en
        \qty{119.9}{pm}. Wijs ze toe en verklaar met behulp van resonantie.
  \item Leg uit waarom de hoek \ce{O-N-O} van \ce{NO2} ($134^\circ$)
        groter is dan $120^\circ$.
  \item Welke van de twee structuren van vraag 4 legt de negatieve
        \omterm{def:b1:lewis-resonance-vsepr:formal-charge}{formele lading} op het elektronegatiefste atoom?
\end{enumerate}

\medskip
\noindent\textbf{Deel III --- Vormen en polariteit.}
\begin{enumerate}[resume]
  \item Geef het type \ce{AX_nE_m} van het centrale atoom van \ce{O3},
        \ce{N2O}, \ce{NO3-} en van het stikstofatoom van \ce{HNO3}, en
        hun vormen.
  \item De hoek van ozon is $116.8^\circ$. Leg uit waarom hij onder
        $120^\circ$ ligt.
  \item Ozon heeft een \omterm{def:b1:lewis-resonance-vsepr:dipole}{dipoolmoment} van \qty{0.53}{\debye}, \ce{N2O}
        \qty{0.16}{\debye}, \ce{N2} geen. Verklaar.
  \item Waarom is \ce{CO2} apolair, terwijl \ce{SO2} polair is?
  \item Voorspel de hoek van \ce{NH3} ten opzichte van $109.5^\circ$ en
        vergelijk daarna met de gemeten $106.7^\circ$.
  \item Dezelfde vraag voor water ($104.5^\circ$).
\end{enumerate}
% ledger: geo:O3.rOO, re:O2, geo:H2O2.rOO, geo:HNO3.rNO, geo:HNO3.rNO2,
% geo:HNO3.rNO3, geo:NO2.aONO, geo:O3.aOOO, dip:O3, dip:N2O

\medskip
\noindent\textbf{Deel IV --- De partiële ladingen van water.}
Het \omterm{def:b1:lewis-resonance-vsepr:dipole}{dipoolmoment} van water is \qty{1.857}{\debye}, zijn hoek
$104.48^\circ$ en de lengte van zijn \ce{O-H}-binding \qty{95.8}{pm}.
% ledger: dip:H2O, geo:H2O.aHOH, geo:H2O.rOH
\begin{enumerate}[resume]
  \item Druk het \omterm{def:b1:lewis-resonance-vsepr:dipole}{dipoolmoment} van water uit als functie van de
        \omterm{def:b1:lewis-resonance-vsepr:dipole}{bindingsdipool} $\mu_{OH}$ en de hoek.
  \item Bereken $\mu_{OH}$ in debye.
  \item Reken het om naar coulombmeter.
  \item Hoe groot zou de \omterm{def:b1:lewis-resonance-vsepr:dipole}{bindingsdipool} zijn als hij bestond uit twee
        volle ladingen $\pm e$ op de twee kernen?
  \item Leid het gedeeltelijk ionaire karakter $\delta$ van de
        \ce{O-H}-binding af: de verhouding van de gemeten \omterm{def:b1:lewis-resonance-vsepr:dipole}{bindingsdipool}
        tot die van volle ladingen.
\end{enumerate}
\end{problem}
